\section{The low cluster of the physical soft matrix and its approximant}

The following auxiliary retains the actual five-region minima and both
restricted residual estimates. It first derives a soft low cluster from
the original gap, then identifies the invariant target of a Hermitian
approximating matrix. The two stages have distinct errors; see
\cite[lines 1195--1255 and 1373--1420]{Arad2016RigorousRG} and
\cite{gap:alvv16_soft_regional_residuals,gap:alvv16_cluster_approximation}.

\begin{lemma}[Physical soft and approximating spectral clusters]
    \label{lem:five_region_soft_and_hermitian_clusters}
    \lean{TNLean.MPS.AreaLaw.ALVV.exists_five_region_soft_clusters}
    \leanok
    Start from the same positive nearest-neighbour interactions,
    five-region minima and simultaneous estimated-center soft matrix
    $G$ of Lemma~\ref{lem:alvv_five_soft_residuals}. Retain all its
    ground data, eigenbases and both residual conclusions. Write
    $C\le4$ for its actual number of crossing bonds and
    $d=10|S|$ for its selected-region correction. The original centered
    matrix has ground energy $E$ and the soft matrix has its attained
    ground energy $\mu$. Suppose the original full ground space $T_0$
    is separated from every other spectral value by $g>0$.

    Let $0<\tau\le1$, set
    $\eta_R=E-\mu+19g/20$ and
    $U_R=2(\eta_R+d+3C+2)$, and suppose
    $T_0'\ge g/40+3C+12$ and $U_R\le t,T_1'$.
    Define
    \begin{align}
        \rho_O&=|S|\left(8C e^{-(T_0'-g/40-3C-12)/8}
                    +\frac{t}{m}(T_0'/t)^m\right),\\
        \rho_R&=|S|\left(8C e^{-(T_1'-U_R)/8}
                    +\frac{t}{m}(T_1'/t)^m\right).
    \end{align}
    If both are at most $\tau g/10^6$, then
    \begin{align}
        T_s=E_G(E+g/20)&=E_G(E+19g/20),
        &\dim T_s&=\dim T_0,
    \end{align}
    and $T_s,T_0$ are mutually $\tau/200$-close. Here
    $E_A(c)$ denotes the closed spectral space of $A$ below $c$.
    The original uncentered ground energy is $E+\sum_j e_j$.

    For any Hermitian $\widehat G$ with $\|G-\widehat G\|\le\xi$,
    $\xi>0$, put $V=\|G\|+\xi$ and
    \begin{align}
        n_-&=\left\lceil\frac{E+g/20-\min\{E+g/20-1,-V\}}{2\xi}\right\rceil,\\
        n_+&=\left\lceil\frac{E+3g/4-\min\{E+3g/4-1,-V\}}{2\xi}\right\rceil,\\
        \delta_-&=\frac{n_-(2\xi)^2}{(g/5+\xi)^2},
        &\delta_+&=\frac{n_+(2\xi)^2}{(g/5+\xi)^2}.
    \end{align}
    If $\delta_-,\delta_+<1$, then
    \begin{align}
        \widehat T=E_{\widehat G}(E+g/4)
          &=E_{\widehat G}(E+3g/4),
        &\dim\widehat T&=\dim T_0.
    \end{align}
    The space $\widehat T$ is $2(\delta_-+\tau/200)$-close to $T_0$,
    while $T_0$ is $2(\delta_++\tau/200)$-close to $\widehat T$.
    The actual minima are chosen before the auxiliary cutoffs.
    No gap or target comparison for either softened matrix is assumed.

    The estimates retain the one-sided ten-unit centers, order $m\ge2$
    and explicit numerical precision conditions. They do not assert a
    uniform parameter choice, measurable adaptive centers or running time.
\end{lemma}
\begin{proof}\leanok
    \uses{lem:alvv_five_soft_residuals}
    Select the physical residual data once. The attained centered ground
    vector identifies the original energy shift and the full ground space
    independently of the eigenbasis. The forward residual is evaluated at
    width $g/40$. The reverse width is exactly $E-\mu+19g/20$, so its
    cutoff is $E+19g/20$. Apply the scaled restricted-residual cluster
    comparison, giving $T_s$ and both errors $\tau/200$.

    Apply low-cluster norm stability to $G,\widehat G$ at the four
    cutoffs $E+g/20,E+g/4,E+3g/4,E+19g/20$. Both inward buffers are
    $g/5$. It gives the approximant's own cluster and the two directed
    losses $\delta_-,\delta_+$. Compose each with the corresponding
    comparison to the original ground space, using equal weights in the
    projection-distance estimate. No residual on a new spectral space
    is inferred from one on a different space.
\end{proof}
